\notechapter{N-20220323}{Complex Differentiability, Cauchy--Riemann Equations, and Holomorphic Rigidity}

\section{One difference quotient must survive every planar direction}

Let \(U\subseteq\mathbb C\) be open and let \(f:U\to\mathbb C\).  The complex derivative at \(z_0\in U\) is
\[
  f'(z_0)
  =\lim_{h\to0}
  \frac{f(z_0+h)-f(z_0)}{h},
\]
provided the limit exists.  The increment \(h\) approaches zero through the whole plane.  This is much stronger than asking for two one-sided real limits.

Write
\[
  z=x+iy,
  \qquad
  f(z)=u(x,y)+iv(x,y).
\]
If \(f\) is complex differentiable at \(z_0=x_0+iy_0\), then the quotient has the same limit along the real and imaginary axes.

Along \(h=t\in\mathbb R\),
\[
  f'(z_0)
  =u_x(x_0,y_0)+iv_x(x_0,y_0).
\]
Along \(h=it\),
\[
\begin{aligned}
  f'(z_0)
  &=\lim_{t\to0}
  \frac{u(x_0,y_0+t)-u(x_0,y_0)
  +i\bigl(v(x_0,y_0+t)-v(x_0,y_0)\bigr)}{it}\\
  &=v_y(x_0,y_0)-iu_y(x_0,y_0).
\end{aligned}
\]
Equality gives the Cauchy--Riemann equations
\[
  \boxed{u_x=v_y,
  \qquad
  u_y=-v_x.}
\]
The derivative can then be written in either form:
\[
  f'=u_x+iv_x=v_y-iu_y.
\]

\section{Cauchy--Riemann is sufficient under real differentiability}

The equations alone at one point do not automatically imply complex differentiability.  A clean sufficient hypothesis is real differentiability.

\begin{proposition}
Suppose \(u,v:\mathbb R^2\to\mathbb R\) are differentiable at \((x_0,y_0)\) and satisfy the Cauchy--Riemann equations there.  Then \(f=u+iv\) is complex differentiable at \(z_0=x_0+iy_0\).
\end{proposition}

\begin{proof}
Let \(h=\Delta x+i\Delta y\).  Real differentiability gives
\[
\begin{aligned}
  \Delta u&=u_x\Delta x+u_y\Delta y+r_1(h),\\
  \Delta v&=v_x\Delta x+v_y\Delta y+r_2(h),
\end{aligned}
\]
where
\[
  \frac{|r_1(h)|+|r_2(h)|}{|h|}\to0.
\]
Using \(u_x=v_y\) and \(u_y=-v_x\),
\[
\begin{aligned}
  \Delta f
  &=\Delta u+i\Delta v\\
  &=(u_x+iv_x)(\Delta x+i\Delta y)
  +r_1(h)+ir_2(h).
\end{aligned}
\]
Hence
\[
  \frac{\Delta f}{h}
  =u_x+iv_x+
  \frac{r_1(h)+ir_2(h)}{h}
  \longrightarrow u_x+iv_x.
\]
\end{proof}

If the first partial derivatives are continuous in a neighborhood, real differentiability is automatic.  Thus the familiar \(C^1\) version is:
\[
  f\text{ holomorphic}
  \quad\Longleftrightarrow\quad
  u_x=v_y,
  \quad
  u_y=-v_x.
\]

\section{Wirtinger derivatives isolate the holomorphic part}

Define
\[
  \partial_z
  =\frac12\left(\partial_x-i\partial_y\right),
  \qquad
  \partial_{\bar z}
  =\frac12\left(\partial_x+i\partial_y\right).
\]
Since
\[
  x=\frac{z+\bar z}{2},
  \qquad
  y=\frac{z-\bar z}{2i},
\]
these operators treat \(z\) and \(\bar z\) as formally independent coordinates:
\[
  \partial_z z=1,
  \quad
  \partial_z\bar z=0,
  \quad
  \partial_{\bar z}z=0,
  \quad
  \partial_{\bar z}\bar z=1.
\]
For \(f=u+iv\),
\[
  2\partial_{\bar z}f
  =(u_x-v_y)+i(v_x+u_y).
\]
Therefore the Cauchy--Riemann equations are exactly
\[
  \boxed{\partial_{\bar z}f=0.}
\]
Under real differentiability, this is equivalent to holomorphicity.

The operator \(\partial_{\bar z}\) measures failure of holomorphicity.  For example,
\[
  \partial_{\bar z}\bar z=1,
\]
so \(\bar z\) is nowhere holomorphic.  Also
\[
  |z|^2=z\bar z,
  \qquad
  \partial_{\bar z}|z|^2=z.
\]
Hence \(|z|^2\) is complex differentiable only at \(z=0\), and it is not holomorphic on any neighborhood of that point.

\section{The Jacobian must be a rotation followed by a scaling}

The real derivative of \(f=u+iv\) is the Jacobian matrix
\[
  Df=
  \begin{pmatrix}
  u_x&u_y\\
  v_x&v_y
  \end{pmatrix}.
\]
Under the Cauchy--Riemann equations,
\[
  Df=
  \begin{pmatrix}
  a&-b\\
  b&a
  \end{pmatrix},
  \qquad
  a=u_x,
  \quad
  b=v_x.
\]
This is precisely the real matrix of multiplication by
\[
  f'(z_0)=a+ib.
\]
If \(f'(z_0)\ne0\), write
\[
  f'(z_0)=\rho e^{i\theta}.
\]
Then
\[
  Df(z_0)=
  \rho
  \begin{pmatrix}
  \cos\theta&-\sin\theta\\
  \sin\theta&\cos\theta
  \end{pmatrix}.
\]
Infinitesimally, \(f\) scales every direction by \(\rho\) and rotates every direction by \(\theta\).  Thus a holomorphic map with nonzero derivative preserves oriented angles.  This is the local meaning of conformality.

At a critical point \(f'(z_0)=0\), the first-order map collapses.  For instance, \(f(z)=z^m\) near \(0\) multiplies angles by \(m\) but has zero linear part when \(m\ge2\).

\section{Cauchy's integral formula turns boundary values into derivatives}

The rigidity of holomorphic functions does not follow from the Cauchy--Riemann equations alone by a short algebraic argument.  The decisive analytic input is Cauchy's integral formula.

If \(f\) is holomorphic on a neighborhood of the closed disk
\[
  \overline{D(a,R)},
\]
then for \(|z-a|<R\),
\[
  \boxed{
  f(z)=\frac{1}{2\pi i}
  \int_{|\zeta-a|=R}
  \frac{f(\zeta)}{\zeta-z}\,d\zeta.}
\]
Differentiating under the integral sign gives
\[
  f^{(n)}(z)
  =\frac{n!}{2\pi i}
  \int_{|\zeta-a|=R}
  \frac{f(\zeta)}{(\zeta-z)^{n+1}}\,d\zeta.
\]
One complex derivative therefore forces derivatives of every order.

The same formula yields the Cauchy estimate.  If
\[
  M_R=\max_{|\zeta-a|=R}|f(\zeta)|,
\]
then
\[
  |f^{(n)}(a)|
  \le \frac{n!M_R}{R^n}.
\]
Boundary control becomes quantitative control of all derivatives in the interior.

\section{Holomorphic functions are analytic}

For \(|z-a|<R\), expand the kernel as a geometric series:
\[
\begin{aligned}
  \frac{1}{\zeta-z}
  &=\frac{1}{\zeta-a}
  \frac{1}{1-\frac{z-a}{\zeta-a}}\\
  &=\sum_{n=0}^{\infty}
  \frac{(z-a)^n}{(\zeta-a)^{n+1}}.
\end{aligned}
\]
The series converges uniformly on the integration circle for \(|z-a|<R\).  Substitution into Cauchy's formula gives
\[
  f(z)=\sum_{n=0}^{\infty}c_n(z-a)^n,
\]
where
\[
  c_n=
  \frac{1}{2\pi i}
  \int_{|\zeta-a|=R}
  \frac{f(\zeta)}{(\zeta-a)^{n+1}}\,d\zeta
  =\frac{f^{(n)}(a)}{n!}.
\]
Thus
\[
  \boxed{
  f(z)=\sum_{n=0}^{\infty}
  \frac{f^{(n)}(a)}{n!}(z-a)^n.}
\]
Holomorphicity implies analyticity.  This is the sharp contrast with smooth real functions such as
\[
  g(x)=
  \begin{cases}
  e^{-1/x^2},&x\ne0,\\
  0,&x=0,
  \end{cases}
\]
whose Taylor series at \(0\) vanishes identically although the function does not.

\section{Zeros are isolated unless the whole function vanishes}

Suppose \(f\) is holomorphic near \(a\) and not identically zero.  Its Taylor series has a first nonzero coefficient:
\[
  f(z)=c_m(z-a)^m+c_{m+1}(z-a)^{m+1}+\cdots,
  \qquad
  c_m\ne0.
\]
Factor
\[
  f(z)=(z-a)^m g(z),
  \qquad
  g(a)=c_m\ne0.
\]
By continuity, \(g\) is nonzero near \(a\), so \(a\) is an isolated zero of multiplicity \(m\).

This proves the identity theorem.  If two holomorphic functions agree on a set with a limit point inside a connected domain, their difference has non-isolated zeros and must vanish identically.

The chain of rigidity is now explicit:
\[
\begin{aligned}
  \partial_{\bar z}f=0
  &\longrightarrow
  \text{Cauchy formula}
  \longrightarrow
  \text{power series},\\
  &\longrightarrow
  \text{isolated zeros and unique continuation}.
\end{aligned}
\]
The operator \(\partial_{\bar z}\) is the local differential equation; Cauchy's formula supplies the global analytic mechanism that makes its solutions rigid.
